\documentclass[a4paper,12pt]{article}

% Пакеты для русского языка и кодировки
\usepackage[utf8]{inputenc}
\usepackage[T2A]{fontenc}
\usepackage[russian]{babel}

% Пакеты для математики
\usepackage{amsmath}
\usepackage{amssymb}
\usepackage{amsthm}

% Пакеты для оформления
\usepackage{geometry}
\geometry{left=2.5cm, right=2.5cm, top=2.5cm, bottom=2.5cm}
\usepackage{hyperref}
\usepackage{enumitem}
\usepackage{booktabs}
\usepackage{array}

% Настройка заголовков
\title{2.3.5 Материалы для изучения по научной специальности \\
Математическое и программное обеспечение вычислительных систем, комплексов и компьютерных сетей \\
Гусев Виталий Евгеньевич 2.3.5 2026}
\date{}

\begin{document}

\maketitle

\tableofcontents
\newpage

\setcounter{section}{1} % чтобы раздел имел номер 2
\section{Раздел 3. Вычислительные машины, системы и сети. Языки программирования.}

\subsection{Классификация вычислительных систем по способу организации параллельной обработки}

\noindent\textit{Основные источники: Орлов С.А., Цилькер Б.Я. Организация ЭВМ и систем: Учебник для вузов. 2-е изд. СПб.: Питер, 2011. Гл. 10, с. 450--465; гл. 11, с. 466--492; гл. 13, с. 526--556; гл. 14, с. 557--578.}

\subsubsection{Понятие параллельной обработки и необходимость классификации}

\paragraph{Ограничения фон-неймановской архитектуры.}
Фон-неймановская архитектура ориентирована на последовательное исполнение команд программы. В условиях постоянно возрастающих требований к производительности вычислительной техники всё очевиднее становятся ограничения такого подхода. Дальнейшее развитие компьютерных средств связано с переходом к параллельным вычислениям — как в рамках одной ВМ, так и путём создания многопроцессорных систем и сетей, объединяющих большое количество отдельных процессоров или отдельных вычислительных машин. Для такого подхода вместо термина «вычислительная машина» более подходит термин «вычислительная система» (ВС). Отличительной особенностью вычислительных систем является наличие в них средств, реализующих параллельную обработку, за счёт построения параллельных ветвей в вычислениях, что не предусматривалось классической структурой ВМ.
\textit{(Источник: Орлов, Цилькер, гл. 10, с. 450.)}

\paragraph{Многообразие параллельных систем и задачи классификации.}
Даже краткое перечисление типов современных параллельных вычислительных систем даёт понять, что для ориентации в этом многообразии необходима чёткая система классификации. От ответа на главный вопрос — что заложить в основу классификации — зависит, насколько конкретная система классификации помогает разобраться с тем, что представляет собой архитектура ВС и насколько успешно данная архитектура позволяет решать определённый круг задач. Попытки систематизировать всё множество архитектур параллельных вычислительных систем предпринимались достаточно давно и длятся по сей день, но к однозначным выводам пока не привели. Исчерпывающий обзор существующих систем классификации ВС приведён в [5]. В данном подразделе рассматриваются наиболее распространённые подходы к классификации: по уровням параллелизма и гранулярности, по классификации Флинна, по способу организации памяти и по способу управления вычислениями.
\textit{(Источник: Орлов, Цилькер, гл. 10, с. 463.)}

\subsubsection{Уровни параллелизма и гранулярность}

\paragraph{Уровни параллелизма.}
Методы и средства реализации параллелизма зависят от того, на каком уровне он должен обеспечиваться. Обычно различают следующие уровни параллелизма:
\begin{itemize}
\item \emph{микроуровень}: выполнение команды разделяется на фазы, а фазы нескольких соседних команд могут быть перекрыты за счёт конвейеризации; уровень достижим на ВС с одним процессором;
\item \emph{уровень команд}: параллельное выполнение нескольких команд достигается посредством размещения в процессоре сразу нескольких конвейеров; реализуется в суперскалярных процессорах;
\item \emph{уровень потоков}: задачи разбиваются на части, которые могут выполняться параллельно (потоки); данный уровень достигается на параллельных ВС;
\item \emph{уровень заданий}: несколько независимых заданий одновременно выполняются на разных процессорах, практически не взаимодействуя друг с другом; этот уровень реализуется на многопроцессорных и многомашинных ВС.
\end{itemize}
\textit{(Источник: Орлов, Цилькер, гл. 10, с. 450--451.)}

\paragraph{Гранулярность параллелизма.}
С понятием уровня параллелизма тесно связано понятие \emph{гранулярности} — меры отношения объёма вычислений, выполненных в параллельной задаче, к объёму коммуникаций (для обмена сообщениями). Степень гранулярности варьируется от мелкозернистой до крупнозернистой. Различают:
\begin{itemize}
\item \emph{крупнозернистый параллелизм} (coarse grained): каждое параллельное вычисление достаточно независимо от остальных, требуется относительно редкий обмен информацией; единицами распараллеливания являются большие и независимые программы, включающие тысячи команд; уровень обеспечивается операционной системой;
\item \emph{среднезернистый параллелизм} (medium grained): единицами распараллеливания являются вызываемые процедуры, включающие сотни команд; обычно организуется как программистом, так и компилятором;
\item \emph{мелкозернистый параллелизм} (fine grained): каждое параллельное вычисление мало и элементарно, составляется из десятков команд; распараллеливаемыми единицами являются элементы выражения или отдельные итерации цикла; характерная особенность — приблизительное равенство интенсивности вычислений и обмена данными; уровень часто используется распараллеливающим (векторизующим) компилятором.
\end{itemize}
Эффективное параллельное исполнение требует искусного баланса между степенью гранулярности программы и величиной коммуникационной задержки. Если коммуникационная задержка минимальна, наилучшую производительность обещает мелкоструктурное разбиение программы (параллелизм данных). Если задержка велика, предпочтительнее крупнозернистое разбиение.
\textit{(Источник: Орлов, Цилькер, гл. 10, с. 451.)}

\paragraph{Профиль параллелизма.}
Число процессоров, параллельно выполняющих программу в каждый момент времени \(t\), задаёт степень параллелизма \(P(t)\) (Degree Of Parallelism). Графическое представление параметра \(P\) в функции времени называют \emph{профилем параллелизма программы}. Изменения числа параллельно работающих процессоров зависят от многих факторов (алгоритма, доступных ресурсов, степени оптимизации, обеспечиваемой компилятором, и т.д.). Типичный профиль параллелизма для алгоритма декомпозиции (divide-and-conquer algorithm) показан на рис. 10.2 (см. учебник). Средний параллелизм \(A\) определяется как
\[
A=\frac{\sum_{i=1}^{n} i t_i}{\sum_{i=1}^{n} t_i},
\]
где \(t_i\) — интервал времени, в течение которого загружены \(i\) процессоров. Профиль параллелизма позволяет оценить, насколько эффективно программа использует доступные процессоры.
\textit{(Источник: Орлов, Цилькер, гл. 10, с. 452--453.)}

\subsubsection{Классификация Флинна}

\paragraph{Понятие потока.}
Среди всех систем классификации ВС наибольшее признание получила классификация, предложенная в 1966 году М. Флинном. В её основу положено понятие потока, под которым понимается последовательность элементов, команд или данных, обрабатываемая процессором. В зависимости от количества потоков команд и потоков данных Флинн выделяет четыре класса архитектур: SISD, MISD, SIMD, MIMD.
\textit{(Источник: Орлов, Цилькер, гл. 10, с. 463.)}

\paragraph{SISD.}
SISD (Single Instruction Stream/Single Data Stream) — одиночный поток команд и одиночный поток данных. Представителями этого класса являются классические фон-неймановские ВМ, где имеется только один поток команд, команды обрабатываются последовательно и каждая команда инициирует одну операцию с одним потоком данных. То, что для увеличения скорости обработки команд и скорости выполнения арифметических операций может применяться конвейерная обработка, не имеет значения, поэтому в класс SISD одновременно попадают как ВМ CDC 6600 со скалярными функциональными устройствами, так и CDC 7600 с конвейерными. Некоторые специалисты считают, что к SISD-системам можно причислить и векторно-конвейерные ВС, если рассматривать вектор как неделимый элемент данных для соответствующей команды.
\textit{(Источник: Орлов, Цилькер, гл. 10, с. 463--464.)}

\paragraph{MISD.}
MISD (Multiple Instruction Stream/Single Data Stream) — множественный поток команд и одиночный поток данных. Из определения следует, что в архитектуре ВС присутствует множество процессоров, обрабатывающих один и тот же поток данных. Примером могла бы служить ВС, на процессоры которой подаётся искажённый сигнал, а каждый из процессоров обрабатывает этот сигнал с помощью своего алгоритма фильтрации. Тем не менее ни Флинн, ни другие специалисты до сих пор не сумели представить убедительный пример реально существующей вычислительной системы, построенной на данном принципе. Ряд исследователей относят к данному классу конвейерные системы, однако это не нашло окончательного признания. Отсюда принято считать, что пока данный класс пуст. Наличие пустого класса не следует считать недостатком классификации; такие классы могут стать полезными для разработки принципиально новых концепций в теории и практике построения ВС.
\textit{(Источник: Орлов, Цилькер, гл. 10, с. 463--464.)}

\paragraph{SIMD.}
SIMD (Single Instruction Stream/Multiple Data Stream) — одиночный поток команд и множественный поток данных. ВМ данной архитектуры позволяют выполнять одну арифметическую операцию сразу над многими данными — элементами вектора. Бесспорными представителями класса SIMD считаются матрицы процессоров, где единое управляющее устройство контролирует множество процессорных элементов. Все процессорные элементы получают от устройства управления одинаковую команду и выполняют её над своими локальными данными. В принципе, в этот класс можно включить и векторно-конвейерные ВС, если каждый элемент вектора рассматривать как отдельный элемент потока данных. К классу SIMD относятся векторные и векторно-конвейерные ВС, матричные ВС, ассоциативные ВС и вычислительные системы с систолической структурой (см. гл. 13).
\textit{(Источник: Орлов, Цилькер, гл. 10, с. 464.)}

\paragraph{MIMD.}
MIMD (Multiple Instruction Stream/Multiple Data Stream) — множественный поток команд и множественный поток данных. Класс предполагает наличие в вычислительной системе множества устройств обработки команд, объединённых в единый комплекс и работающих каждое со своим потоком команд и данных. Класс MIMD чрезвычайно широк, поскольку включает в себя всевозможные мультипроцессорные системы. Кроме того, приобщение к классу MIMD зависит от трактовки: так, векторно-конвейерные ВС можно вполне отнести и к классу MIMD, если конвейерную обработку рассматривать как выполнение множества команд (операций ступеней конвейера) над множественным скалярным потоком. К классу MIMD относятся симметричные мультипроцессорные системы (SMP), системы с неоднородным доступом к памяти (NUMA), системы с массовой параллельной обработкой (MPP), кластеры и Constellation-системы (см. гл. 14).
\textit{(Источник: Орлов, Цилькер, гл. 10, с. 464--465.)}

Схема классификации Флинна вплоть до настоящего времени является наиболее распространённой при первоначальной оценке той или иной ВС, поскольку позволяет сразу оценить базовый принцип работы системы. Однако у неё имеются и очевидные недостатки, например неспособность однозначно отнести некоторые архитектуры к тому или иному классу, а также чрезмерная насыщенность класса MIMD. Это породило множественные попытки либо модифицировать классификацию Флинна, либо предложить иную систему классификации.
\textit{(Источник: Орлов, Цилькер, гл. 10, с. 465.)}

\subsubsection{Классификация по способу организации памяти}

\paragraph{Разделяемая и распределённая память.}
В вычислительных системах, объединяющих множество параллельно работающих процессоров или машин, задача правильной организации памяти является одной из важнейших. В зависимости от того, каким образом организована память многопроцессорных (многомашинных) систем, различают вычислительные системы с разделяемой памятью (shared memory) и ВС с распределённой памятью (distributed memory). В системах с разделяемой памятью память ВС рассматривается как общий ресурс, и каждый из процессоров имеет полный доступ ко всему адресному пространству. Такие системы называют сильно связанными (closely coupled systems). В варианте с распределённой памятью каждому из процессоров придаётся собственная память; процессоры объединяются в сеть и могут при необходимости обмениваться данными, хранящимися в их памяти, передавая друг другу сообщения. Такой вид ВС называют слабо связанными (loosely coupled systems).
\textit{(Источник: Орлов, Цилькер, гл. 11, с. 466--467.)}

\paragraph{Различие в интерпретации адреса.}
Различие между разделяемой и распределённой памятью заключается в структуре виртуальной памяти, то есть в том, как память выглядит со стороны процессора. Физически почти каждая система памяти разделена на автономные компоненты, доступ к которым может производиться независимо. Разделяемую память от распределённой отличает то, каким образом подсистема памяти интерпретирует поступивший от процессора адрес ячейки. В случае разделяемой памяти \(i\) — это глобальный адрес, и для любого процессора он указывает на одну и ту же ячейку. В распределённой системе памяти \(i\) — это локальный адрес: если два процессора выполняют команду load \(R0, i\), то каждый из них обращается к \(i\)-й ячейке в своей локальной памяти, то есть к разным ячейкам, и в регистры \(R0\) могут быть загружены неодинаковые значения. Различие между двумя системами памяти должно учитываться программистом, поскольку оно определяет способ взаимодействия частей распараллеленной программы. В варианте с разделяемой памятью достаточно создать в памяти структуру данных и передавать в параллельно используемые подпрограммы ссылки на эту структуру. В системе с распределённой памятью необходимо в каждой локальной памяти иметь копию совместно используемых данных; эти копии создаются путём вкладывания разделяемых данных в сообщения, посылаемые другим процессорам.
\textit{(Источник: Орлов, Цилькер, гл. 11, с. 467.)}

\paragraph{Модели архитектур памяти.}
В рамках как разделяемой, так и распределённой памяти реализуются несколько архитектур системы памяти. На рис. 11.1 (см. учебник) приведена классификация таких моделей. Для систем с физически разделяемой памятью характерна модель UMA (Uniform Memory Access) — однородный доступ к памяти: все процессоры имеют равные возможности по доступу к единому адресному пространству, доступ любого процессора к памяти производится единообразно и занимает одинаковое время. Единая память может быть построена как одноблочная или по блочному принципу, но обычно практикуется второй вариант. UMA-системы предполагают наличие узла, соединяющего каждый из \(n\) процессоров с каждым из \(m\) банков памяти. Простейший путь построения таких ВС — объединение нескольких процессоров с банками разделяемой памяти посредством общей шины; альтернативный способ — использование коммутатора, маршрутизирующего запросы процессора к одному из нескольких банков памяти.
\textit{(Источник: Орлов, Цилькер, гл. 11, с. 467--469.)}

Для систем с физически распределённой разделяемой памятью характерны модели NUMA (Non-Uniform Memory Access), COMA (Cache Only Memory Architecture) и DSM (Distribute Shared Memory). В NUMA-системах каждый процессор имеет локальную память, которая рассматривается как часть общей памяти ВС и которой выделен отдельный диапазон адресов в едином адресном пространстве. Помимо возможности обратиться к своей локальной памяти, каждый процессор имеет возможность физически «добраться» до локальной памяти остальных процессоров. Неоднородность обращения к памяти выражается в том, что доступ процессора к собственной локальной памяти производится напрямую (это намного быстрее, чем доступ к удалённой памяти через коммутатор или сеть). В COMA локальная память каждого процессора считается частью большой кэш-памяти; совокупность кэш-памяти всех процессоров рассматривается как глобальная память системы, причём собственно глобальная память отсутствует. В DSM каждый процессор обладает локальной памятью, но в отличие от NUMA не имеет возможности физического доступа к локальной памяти других процессоров; ВС представляется пользователю как система с общей памятью лишь благодаря операционной системе, которая предлагает единое адресное пространство, а фактическое обращение к памяти «чужого» компьютера обеспечивается путём обмена сообщениями между процессорами.
\textit{(Источник: Орлов, Цилькер, гл. 11, с. 470--472.)}

Для систем с распределённой памятью характерна модель NORMA (No Remote Memory Access): каждый процессор имеет доступ только к своей локальной памяти; доступ к удалённой памяти возможен только путём обмена сообщениями с процессором, которому принадлежит адресуемая память. Подобная организация имеет ряд достоинств: при доступе к данным не возникает конкуренции за шину или коммутаторы; размер системы ограничивает только сеть, объединяющая процессоры; снимается проблема когерентности кэш-памяти. Основной недостаток — сложность обмена информацией между процессорами, что приводит к двум видам издержек: требуется время на формирование и пересылку сообщения; для обеспечения реакции на сообщения принимающий процессор должен получить запрос прерывания и выполнить процедуру обработки этого прерывания.
\textit{(Источник: Орлов, Цилькер, гл. 11, с. 472--473.)}

\subsubsection{Классификация MIMD-систем}

\paragraph{Общая характеристика.}
В MIMD-системе каждый процессорный элемент (ПЭ) выполняет свою программу достаточно независимо от других ПЭ. В то же время ПЭ должны как-то взаимодействовать друг с другом. Различие в способе такого взаимодействия определяет условное деление MIMD-систем на ВС с разделяемой памятью и системы с распределённой памятью. В системах с разделяемой памятью, которые характеризуют как сильно связанные (tightly coupled), реализовано единое адресное пространство, доступное всем процессорным элементам посредством общей шины или сети соединений. Такие ВС часто называют мультипроцессорами. В системах с распределённой памятью или слабо связанных (loosely coupled) вычислительных системах вся память распределена между процессорными элементами, и каждый блок памяти доступен только «своему» процессору. Фактически в качестве элементов ВС выступают не процессоры, а вычислительные машины, связанные между собой сетью соединений. Иногда для обозначения подобных ВС используют термин «мультикомпьютеры».
\textit{(Источник: Орлов, Цилькер, гл. 14, с. 557.)}

\paragraph{Симметричные мультипроцессорные системы (SMP).}
SMP-система состоит из множества процессорных элементов, каждый из которых имеет равноправный доступ к логически единой памяти (физически память обычно строится по блочному принципу). Таким образом, в системе реализована концепция однородного доступа к памяти (UMA). Обычно все ПЭ идентичны, однако допускается использование и различных ПЭ при условии сопоставимой производительности. Все процессорные элементы управляются единственным экземпляром операционной системы, загружаемой в совместно используемую память. Типовая SMP-система содержит от двух до 32 идентичных процессоров, в качестве которых обычно выступают недорогие RISC-процессоры; в последнее время наметилась тенденция оснащения SMP-систем также и CISC-процессорами. Каждый процессор снабжён локальной кэш-памятью; согласованность содержимого кэш-памяти всех процессоров обеспечивается аппаратными средствами (протоколы когерентности). Все процессоры ВС имеют равноправный доступ к разделяемой основной памяти и устройствам ввода/вывода. Важным аспектом архитектуры симметричных мультипроцессоров является способ взаимодействия процессоров с общими ресурсами; с этих позиций выделяют архитектуры с общей шиной, с коммутатором типа «кроссбар», с многопортовой памятью и с централизованным устройством управления. На практике наиболее распространена архитектура с общей шиной. Основной недостаток SMP-систем — при увеличении числа процессоров стоимость системы растёт быстрее, чем производительность, а из-за задержек при одновременном обращении нескольких ПЭ к общей памяти происходит взаимное торможение.
\textit{(Источник: Орлов, Цилькер, гл. 14, с. 558--562.)}

\paragraph{Системы с неоднородным доступом к памяти (NUMA).}
В отличие от SMP-систем с разделяемой памятью, где максимальное количество ПЭ колеблется в диапазоне 16--64, в системах с неоднородным доступом к памяти (NUMA, Non-Uniform Memory Access) реализована технология, позволяющая создавать крупномасштабные вычислительные системы. В архитектуре NUMA память физически распределена, но логически общедоступна. Имеется множество независимых составляющих ВС (узлов), объединённых с помощью какой-либо сети соединений (например, кроссбара, кольца и т.д.). Узел содержит процессор с кэш-памятью, а также локальную основную память, рассматриваемую как часть глобальной основной памяти системы. Каждый узел владеет локальной памятью, но с позиций системы имеет место глобальное адресное пространство, где каждая ячейка любой локальной основной памяти имеет уникальный системный адрес. Когда процессор инициирует доступ к памяти и нужная ячейка отсутствует в его локальной кэш-памяти, организуется операция выборки: если нужная ячейка находится в локальной памяти, выборка производится с использованием локальной шины; если требуемая ячейка хранится в удалённой секции глобальной памяти, автоматически формируется запрос, посылаемый по сети соединений на локальную шину узла, где находится запрошенная информация. В подавляющем большинстве NUMA-систем реализована аппаратная поддержка когерентности кэш-памяти процессорных элементов (ccNUMA), хотя известны системы, где такая поддержка отсутствует (nccNUMA). Для взаимодействия узлов используется не шина, а сеть соединений с более сложной топологией; в ccNUMA-системах проблема когерентности решается с помощью протоколов на основе распределённых справочников. Масштабируемость NUMA-систем ограничена лишь величиной адресного пространства, возможностями аппаратных средств поддержки когерентности кэшей и возможностями операционной системы по управлению большим числом процессоров. Благодаря своей архитектуре и заложенному в неё потенциалу масштабируемости NUMA-системы считаются перспективным направлением при создании мощных вычислительных систем.
\textit{(Источник: Орлов, Цилькер, гл. 14, с. 563--565.)}

\paragraph{Системы с массовой параллельной обработкой (MPP).}
Основные причины появления систем с массовой параллельной обработкой (MPP, Massively Parallel Processing) — необходимость построения ВС с гигантской производительностью и стремление раздвинуть границы производства ВС в большом диапазоне производительности и стоимости. MPP-система состоит из множества однородных вычислительных узлов, число которых может исчисляться тысячами. Узел содержит полный комплекс устройств, необходимых для независимого функционирования (процессор, память, подсистему ввода/вывода, коммуникационное оборудование), то есть, по сути, является полноценной вычислительной машиной. Узлы объединены коммуникационной сетью с высокой пропускной способностью и малыми задержками. Работу узлов координирует главная управляющая вычислительная машина (хост-компьютер). Это может быть отдельная ВМ или один из узлов системы. Если система содержит выделенный хост-компьютер, то полноценная ОС функционирует только на нём, а на узлы устанавливается её урезанный вариант, поддерживающий лишь функции ядра ОС; при отсутствии главной ВМ полноценная ОС устанавливается на каждый узел. Процессоры в узлах имеют доступ только к своей локальной памяти; для доступа к памяти другого узла пара узлов должна обменяться сообщениями. Такой режим исключает возможность конфликтов, возникающих в разделяемой памяти, и снимает проблему когерентности кэш-памяти, что позволяет наращивать количество процессоров до нескольких тысяч. Слабым местом MPP является хост-компьютер — при выходе его из строя вся система оказывается неработоспособной. Несмотря на сложности, сфера применения ВС с массовым параллелизмом постоянно расширяется; MPP-системы являются сегодня лидерами по достигнутой производительности.
\textit{(Источник: Орлов, Цилькер, гл. 14, с. 566--568.)}

\paragraph{Кластерные вычислительные системы.}
Кластером называют группу взаимно соединённых вычислительных систем (узлов), работающих совместно и составляющих единый вычислительный ресурс, создавая иллюзию наличия единственной ВМ. Изначально перед кластерами ставились две задачи: достичь большой вычислительной мощности и обеспечить повышенную надёжность ВС. Архитектура кластерных систем во многом похожа на архитектуру MPP-систем: тот же принцип распределённой памяти, использование в качестве вычислительных узлов законченных вычислительных машин, большой потенциал для масштабирования. В качестве узла кластера может выступать как однопроцессорная ВМ, так и ВС типа SMP. Как правило, это не специализированные устройства, как в MPP, а серийно выпускаемые вычислительные машины и системы. Кластерные системы с одинаковыми узлами называют гомогенными, а с разнотипными узлами — гетерогенными. Использование машин массового производства существенно снижает стоимость ВС, а возможность варьирования различных по типу узлов позволяет получить необходимую производительность за приемлемую цену. Узлы в кластерной системе объединены высокоскоростной сетью. Неотъемлемая часть кластера — специализированное программное обеспечение, организующее бесперебойную работу при отказе одного или нескольких узлов. Кластеры обеспечивают высокий уровень доступности — в них отсутствуют единая операционная система и совместно используемая память, то есть нет проблемы когерентности кэшей. Кластеры хорошо масштабируются путём добавления узлов, что позволяет достичь высочайших показателей производительности.
\textit{(Источник: Орлов, Цилькер, гл. 14, с. 569--572.)}

\paragraph{Кластеры больших SMP-систем (Constellations).}
Огромный потенциал масштабирования, свойственный кластерной архитектуре, делает её очень перспективным направлением в области создания высокопроизводительных вычислительных систем. Масштабирование возможно как за счёт увеличения числа узлов, так и путём применения в качестве узлов не одиночных ВМ, а хорошо масштабируемых вычислительных систем, обычно SMP-типа. Это направление получило настолько широкое развитие, что было выделено в отдельную группу MIMD-устройств — Constellations (термин, обусловленный названием одной из первых ВС данного типа — Sun «Constellation», можно перевести как «созвездие»). Constellation-система — это кластер, узлами которого служат SMP-системы. В качестве отличительного признака выступает число процессорных элементов в узле: в настоящее время условием причисления ВС к Constellation-системам служит число ПЭ в узле — оно должно быть больше (равно) 16; системы, не отвечающие данному условию, считаются классическими кластерными системами. Перспективность Constellation-систем обусловлена удачным сочетанием преимуществ распределённой памяти (возможности наращивания количества узлов) и разделяемой памяти (эффективного доступа множества ПЭ к памяти). Формально структура Constellation-системы полностью соответствует кластерной архитектуре, однако явная направленность на высокую производительность может отражаться в некоторых конструктивных решениях.
\textit{(Источник: Орлов, Цилькер, гл. 14, с. 572--573.)}

\subsubsection{Классификация по способу управления вычислениями}

\paragraph{Управление последовательностью команд.}
Традиционные (фон-неймановские) вычислительные системы, управляемые с помощью счётчика команд, иногда называют вычислительными системами, управляемыми последовательностью команд (control flow computer). Согласно этой концепции порядок вычислений определяет программа, представленная последовательностью команд; команды выполняются в порядке их следования в программе. Это обеспечивается наличием в каждом процессоре счётчика команд. Выполнение команд в каждом процессоре — поочерёдное и потому достаточно медленное. Для получения выигрыша программист или компилятор должны определить независимые команды, которые могут быть поданы на отдельные процессоры.
\textit{(Источник: Орлов, Цилькер, гл. 15, с. 579.)}

\paragraph{Управление от потока данных.}
В вычислительных системах с управлением от потока данных (dataflow) команда выполняется, когда становятся доступными её операнды. Вычислительная модель описывается ориентированным графом потоков данных, вершины которого соответствуют операциям, а дуги — потокам данных (токенам). Операция в вершине выполняется лишь когда по дугам в неё поступила вся необходимая информация. Одновременно могут быть активированы несколько вершин, при этом параллелизм в вычислительной модели выявляется автоматически. Статические потоковые системы допускают не более одного токена на дуге; динамические системы (с помеченными или явно адресуемыми токенами) снимают это ограничение за счёт учёта контекста (цвета) токенов. Мультипотоковая обработка представляет собой крупнозернистую потоковую обработку, в которой вершины графа соответствуют последовательностям команд (нитям).
\textit{(Источник: Орлов, Цилькер, гл. 15, с. 580--591.)}

\paragraph{Управление по запросу.}
В системах с управлением по запросу (demand-driven, редукционных) вычисления инициируются не по готовности данных, а на основе запроса на данные. Процесс представляет собой редукцию графа выражения: выбираются и упрощаются те части выражения, которые можно свести к более простому виду. Вычисления являются «ленивыми»: операция выполняется только тогда, когда её результат требуется другой операции. Различают строчную и графовую редукционные модели. Математической основой редукционных ВС служит лямбда-исчисление; программирование ведётся на функциональных языках. Основные преимущества потоковых и редукционных моделей — естественный параллелизм и отсутствие счётчика команд; основные недостатки — накладные расходы на пересылку токенов, согласование и потерю локальности.
\textit{(Источник: Орлов, Цилькер, гл. 15, с. 593--596.)}

\subsubsection{Метрики параллельных вычислений}

\paragraph{Ускорение и эффективность.}
Метрики параллельных вычислений — это система показателей, позволяющая оценить преимущества, получаемые при параллельном решении задачи на \(n\) процессорах, по сравнению с последовательным решением той же задачи на единственном процессоре. Базисом для определения метрик являются следующие характеристики вычислений: \(n\) — количество процессоров; \(O(n)\) — объём вычислений, выраженный через количество операций, выполняемых \(n\) процессорами; \(T(n)\) — общее время вычислений с использованием \(n\) процессоров. Ускорение (Speedup) за счёт параллельного выполнения программы служит показателем эффективной скорости вычислений:
\[
S(n)=\frac{T(1)}{T(n)}.
\]
Эффективность (Efficiency) характеризует целесообразность наращивания числа процессоров через ту долю ускорения, достигнутого за счёт параллельных вычислений, которая приходится на один процессор:
\[
E(n)=\frac{S(n)}{n}=\frac{T(1)}{nT(n)}.
\]
Идеальное ускорение равно \(n\), идеальная эффективность равна 1.
\textit{(Источник: Орлов, Цилькер, гл. 10, с. 452--454.)}

\paragraph{Закон Амдала.}
Закон Амдала (1967) отражает зависимость ускорения вычислений от числа процессоров и от соотношения между последовательной и распараллеливаемой частями программы при условии, что объём решаемой задачи остаётся неизменным. Если доля операций, которые должны выполняться последовательно одним из процессоров, равна \(f\) (\(0\le f\le1\)), то ускорение на \(n\) процессорах равно
\[
S(n)=\frac{1}{f+\frac{1-f}{n}}.
\]
При безграничном увеличении числа процессоров ускорение стремится к \(1/f\). Это означает, что если в программе 25\% последовательных операций (\(f=0{,}25\)), то сколько бы процессоров ни использовалось, ускорения более чем в четыре раза получить нельзя. Закон Амдала служит аргументом против вычислительных систем с очень большим числом процессоров, решающих общую задачу, поскольку издержки на обмен информацией между процессорами могут оказаться существенными.
\textit{(Источник: Орлов, Цилькер, гл. 10, с. 457--459.)}

\paragraph{Закон Густафсона.}
Закон Густафсона (1988) исходит из того, что при увеличении числа процессоров пользователь обычно сохраняет время вычислений, но пропорционально увеличивает объём решаемой задачи. В этом случае ускорение является линейной функцией числа процессоров:
\[
S(n)=f+(1-f)n.
\]
Закон Густафсона не противоречит закону Амдала: различие состоит лишь в форме использования дополнительной мощности ВС, возникающей вследствие увеличения числа процессоров. Закон Сана–Ная (1993) обобщает оба закона, учитывая ограничение на объём задачи, связанное с ёмкостью доступной памяти.
\textit{(Источник: Орлов, Цилькер, гл. 10, с. 459--462.)}

\subsubsection{Краткое резюме подраздела}
Классификация вычислительных систем по способу организации параллельной обработки включает несколько взаимодополняющих подходов. По уровням параллелизма различают микроуровень (конвейеризация), уровень команд (суперскалярность), уровень потоков и уровень заданий; по гранулярности — крупно-, средне- и мелкозернистый параллелизм. Классификация Флинна выделяет четыре класса архитектур: SISD, MISD, SIMD, MIMD; класс MISD остаётся пустым. По способу организации памяти различают системы с разделяемой памятью (сильно связанные, UMA, NUMA, COMA, DSM) и системы с распределённой памятью (слабо связанные, NORMA). Класс MIMD включает SMP-системы, NUMA-системы, MPP-системы, кластеры и Constellation-системы. По способу управления вычислениями различают системы с управлением последовательностью команд (control flow), управлением от потока данных (dataflow) и управлением по запросу (редукционные). Метрики параллельных вычислений (ускорение, эффективность) и законы Амдала, Густафсона, Сана–Ная позволяют количественно оценить эффективность параллельной обработки и выбрать рациональное число процессоров.
\textit{(Обобщение по: Орлов, Цилькер, гл. 10, с. 450--465; гл. 11, с. 466--492; гл. 13, с. 526--556; гл. 14, с. 557--578; гл. 15, с. 579--597.)}

\subsection{Многопроцессорные и многомашинные комплексы}

\noindent\textit{Основные источники: Орлов С.А., Цилькер Б.Я. Организация ЭВМ и систем: Учебник для вузов. 2-е изд. СПб.: Питер, 2011. Гл. 11, с. 466--492; гл. 14, с. 557--578.}

\subsubsection{Общая характеристика и классификация}

\paragraph{Понятие многопроцессорной и многомашинной системы.}
Под \emph{многопроцессорной вычислительной системой} понимается комплекс, содержащий два или более процессоров (или процессорных элементов), которые совместно используют общую память и работают под управлением единой операционной системы. В отличие от неё, \emph{многомашинная система} (многомашинный комплекс) представляет собой совокупность отдельных вычислительных машин, объединённых коммуникационной сетью для совместного решения задач; каждая машина сохраняет самостоятельность и может работать под управлением собственной операционной системы. В учебнике Орлова и Цилькера многопроцессорные и многомашинные комплексы рассматриваются в рамках класса MIMD (Multiple Instruction Stream/Multiple Data Stream) и классифицируются по способу организации памяти: системы с разделяемой памятью (сильно связанные) и системы с распределённой памятью (слабо связанные).
\textit{(Источник: Орлов, Цилькер, гл. 11, с. 466--467; гл. 14, с. 557.)}

\paragraph{Сильно и слабо связанные системы.}
В системах с разделяемой памятью (shared memory) память ВС рассматривается как общий ресурс, и каждый из процессоров имеет полный доступ ко всему адресному пространству. Такие системы называют \emph{сильно связанными} (closely coupled systems). В варианте с распределённой памятью (distributed memory) каждому из процессоров придаётся собственная память; процессоры объединяются в сеть и могут при необходимости обмениваться данными, хранящимися в их памяти, передавая друг другу сообщения. Такой вид ВС называют \emph{слабо связанными} (loosely coupled systems). Различие между разделяемой и распределённой памятью заключается в структуре виртуальной памяти, то есть в том, как память выглядит со стороны процессора. В системе с разделяемой памятью адрес \(i\) для любого процессора указывает на одну и ту же ячейку; в распределённой системе адрес \(i\) является локальным для каждого процессора. Различие между двумя системами памяти должно учитываться программистом, поскольку оно определяет способ взаимодействия частей распараллеленной программы.
\textit{(Источник: Орлов, Цилькер, гл. 11, с. 466--467.)}

\paragraph{Классификация MIMD-систем.}
В рамках MIMD-систем можно выделить три семейства с распределённой памятью: системы с массовой параллельной обработкой (MPP), кластеры вычислительных машин и кластеры больших SMP-систем (Constellations). Системы с разделяемой памятью включают симметричные мультипроцессорные системы (SMP), параллельные векторные системы (PVP) и системы с неоднородным доступом к памяти (NUMA). Каждая из этих архитектур имеет свои особенности, достоинства и недостатки, определяющие область её применения.
\textit{(Источник: Орлов, Цилькер, гл. 14, с. 558--573.)}

\subsubsection{Симметричные мультипроцессорные системы (SMP)}

\paragraph{Определение и архитектура.}
SMP-система состоит из множества процессорных элементов, каждый из которых имеет равноправный доступ к логически единой памяти (физически память обычно строится по блочному принципу). Таким образом, в системе реализована концепция однородного доступа к памяти (UMA, Uniform Memory Access). Обычно все ПЭ идентичны, однако допускается использование и различных ПЭ при условии сопоставимой производительности. Все процессорные элементы управляются единственным экземпляром операционной системы, загружаемой в совместно используемую память. ОС планирует распределение выполняемых заданий между процессорными элементами: выделяет в процессах фрагменты (нити) и образует из этих фрагментов единую очередь; при освобождении какого-либо ПЭ ему сразу же передаётся очередной фрагмент из очереди.
\textit{(Источник: Орлов, Цилькер, гл. 14, с. 558--559.)}

\paragraph{Характеристики и примеры.}
Типовая SMP-система содержит от двух до 32 идентичных процессоров, в качестве которых обычно выступают недорогие RISC-процессоры. В последнее время наметилась тенденция оснащения SMP-систем также и CISC-процессорами. Каждый процессор снабжён локальной кэш-памятью; согласованность содержимого кэш-памяти всех процессоров обеспечивается аппаратными средствами (протоколы когерентности, например MESI). Все процессоры ВС имеют равноправный доступ к разделяемой основной памяти и устройствам ввода/вывода. Важным аспектом архитектуры симметричных мультипроцессоров является способ взаимодействия процессоров с общими ресурсами; с этих позиций выделяют архитектуры с общей шиной, с коммутатором типа «кроссбар», с многопортовой памятью и с централизованным устройством управления. На практике наиболее распространена архитектура с общей шиной. Примеры SMP-систем: Compaq AlphaServer GS140 и 8400 (до 14 процессоров Alpha 21264), HP N9000 (до 8 процессоров PA-8500), IBM RS/6000 (2–4 процессора PowerPC 604), Sun Enterprise 10000 (64 процессора, коммутатор Gigaplane-XB).
\textit{(Источник: Орлов, Цилькер, гл. 14, с. 559--561.)}

\paragraph{Достоинства и недостатки.}
SMP-система представляется конечному пользователю как единая ВМ высокой производительности. В этом качестве небольшие SMP-системы из 2–4 процессорных элементов часто используются в качестве процессорных элементов в MIMD-системах других видов. Негативные свойства SMP-систем: во-первых, при увеличении числа процессоров стоимость системы растёт быстрее, чем производительность; во-вторых, из-за задержек при одновременном обращении нескольких ПЭ к общей памяти (даже если они параллельно выполняют независимые программы) происходит взаимное торможение.
\textit{(Источник: Орлов, Цилькер, гл. 14, с. 562.)}

\subsubsection{Системы с неоднородным доступом к памяти (NUMA)}

\paragraph{Архитектура NUMA.}
В отличие от SMP-систем с разделяемой памятью, где максимальное количество ПЭ колеблется в диапазоне 16–64, в системах с неоднородным доступом к памяти (NUMA, Non-Uniform Memory Access) реализована технология, позволяющая создавать крупномасштабные вычислительные системы. В архитектуре NUMA память физически распределена, но логически общедоступна. Это позволяет сохранить преимущества архитектуры с единым адресным пространством, а также ощутимо расширяет возможности масштабирования ВС. Имеется множество независимых составляющих ВС (узлов), объединённых с помощью какой-либо сети соединений (например, кроссбара, кольца и т.д.). Узел содержит процессор с кэш-памятью, а также локальную основную память, рассматриваемую как часть глобальной основной памяти системы. Каждый узел владеет локальной памятью, но с позиций системы имеет место глобальное адресное пространство, где каждая ячейка любой локальной основной памяти имеет уникальный системный адрес. Когда процессор инициирует доступ к памяти и нужная ячейка отсутствует в его локальной кэш-памяти, организуется операция выборки: если нужная ячейка находится в локальной памяти, выборка производится с использованием локальной шины; если требуемая ячейка хранится в удалённой секции глобальной памяти, автоматически формируется запрос, посылаемый по сети соединений на локальную шину узла, где находится запрошенная информация. Все эти действия выполняются автоматически, прозрачны для процессора и его кэш-памяти.
\textit{(Источник: Орлов, Цилькер, гл. 14, с. 563--564.)}

\paragraph{Когерентность кэшей в NUMA.}
Как и в любой ВС с разделяемой памятью, особое внимание уделяется когерентности кэшей. В подавляющем большинстве NUMA-систем реализована аппаратная поддержка когерентности кэш-памяти процессорных элементов (ccNUMA), хотя известны системы, где такая поддержка отсутствует (nccNUMA). Поскольку для взаимодействия узлов системы используется не шина, а сеть соединений с более сложной топологией, в ccNUMA-системах проблема когерентности решается с помощью протоколов на основе распределённых справочников. Хотя отдельные реализации и отличаются в деталях, общим является то, что каждый узел содержит справочник; взаимодействуя между собой, справочники позволяют определить физическое расположение любой информации в глобальном адресном пространстве. В nccNUMA согласованность глобальных данных на аппаратном уровне не обеспечивается; забота об использовании таких данных полностью возлагается на программное обеспечение (приложения или компиляторы). Несмотря на этот недостаток, архитектура значительно упрощает масштабирование ВС, то есть увеличение числа процессоров.
\textit{(Источник: Орлов, Цилькер, гл. 14, с. 564--565.)}

\paragraph{Масштабируемость и примеры.}
Масштабируемость NUMA-систем ограничена лишь величиной адресного пространства, возможностями аппаратных средств поддержки когерентности кэшей и возможностями операционной системы по управлению большим числом процессоров. Например, NUMA-система Silicon Graphics Origin поддерживает до 1024 процессоров R10000, а система Sequent NUMA-Q объединяет 252 процессора Pentium II. Очередным этапом развития технологии NUMA стала архитектура NumaFlex, используемая в семействе SGI 3000, где допускается наращивание системы даже за счёт различных процессоров. Благодаря своей архитектуре и заложенному в неё потенциалу масштабируемости NUMA-системы считаются перспективным направлением при создании мощных вычислительных систем.
\textit{(Источник: Орлов, Цилькер, гл. 14, с. 565.)}

\subsubsection{Системы с массовой параллельной обработкой (MPP)}

\paragraph{Назначение и структура.}
Основные причины появления систем с массовой параллельной обработкой (MPP, Massively Parallel Processing) — необходимость построения ВС с гигантской производительностью и стремление раздвинуть границы производства ВС в большом диапазоне производительности и стоимости. MPP-система состоит из множества однородных вычислительных узлов, число которых может исчисляться тысячами. Узел содержит полный комплекс устройств, необходимых для независимого функционирования (процессор, память, подсистему ввода/вывода, коммуникационное оборудование), то есть, по сути, является полноценной вычислительной машиной. Узлы объединены коммуникационной сетью с высокой пропускной способностью и малыми задержками.
\textit{(Источник: Орлов, Цилькер, гл. 14, с. 566.)}

\paragraph{Управление и особенности.}
Работу узлов MPP-системы координирует главная управляющая вычислительная машина (хост-компьютер). Это может быть отдельная ВМ, как, например, в Cray T3D, или один из узлов системы; в последнем случае функции главной ВМ возлагаются на какой-то определённый узел ВС (в течение всего сеанса вычислений). Напомним, что в SMP-системе особые полномочия одному из ПЭ придаются только на время загрузки системы, после чего он опять становится равноправным элементом системы. Если система содержит выделенный хост-компьютер, то полноценная операционная система функционирует только на нём, а на узлы устанавливается её урезанный вариант, поддерживающий лишь функции ядра ОС. При отсутствии главной ВМ полноценная ОС устанавливается на каждый узел MPP-системы; таким образом, каждый узел функционирует под управлением собственной операционной системы. Хост-компьютер (или его заменитель из числа узлов) распределяет задания между множеством подчинённых ему узлов: формирует очередь заданий, каждому из которых назначается некоторый уровень приоритета; по мере освобождения узлов им передаются задания из очереди; узлы оповещают хост-компьютер о ходе выполнения задания, в частности о завершении выполнения или о потребности в дополнительных ресурсах; у хост-компьютера имеются средства для контроля работы подчинённых процессоров, в том числе для обнаружения нештатных ситуаций и прерывания выполнения задания в случае появления более приоритетной задачи.
\textit{(Источник: Орлов, Цилькер, гл. 14, с. 566--567.)}

\paragraph{Преимущества и ограничения.}
Процессоры в узлах имеют доступ только к своей локальной памяти. Для доступа к памяти другого узла пара узлов (клиент и «сервер») должна обменяться сообщениями. Такой режим исключает возможность конфликтов, возникающих в разделяемой памяти при одновременном обращении к ней со стороны нескольких процессоров, и снимает проблему когерентности кэш-памяти. Как следствие, появляется возможность наращивания количества процессоров до нескольких тысяч. Основным признаком, по которому вычислительную систему относят к MPP-типу, часто служит количество процессоров \(n\); строгой границы не существует, но обычно при \(n\ge128\) считается, что это уже MPP. Отличительной чертой MPP является архитектура, позволяющая добиться высочайшей производительности. Благодаря свойству масштабируемости MPP-системы являются сегодня лидерами по достигнутой производительности. Так, список наиболее производительных ВС на конец 2009 года возглавляла MPP-система Jaguar Cray XT5-HE с теоретической пиковой производительностью 2331 TFLOPS; система состоит из 224\,256 вычислительных ядер. Слабым местом MPP является хост-компьютер — при выходе его из строя вся система оказывается неработоспособной. Повышение надёжности главной ВМ лежит на путях упрощения аппаратуры хост-компьютера и/или её дублирования.
\textit{(Источник: Орлов, Цилькер, гл. 14, с. 567--568.)}

\subsubsection{Многомашинные комплексы и кластеры}

\paragraph{Кластерные вычислительные системы.}
Кластером называют группу взаимно соединённых вычислительных систем (узлов), работающих совместно и составляющих единый вычислительный ресурс, создавая иллюзию наличия единственной ВМ. Изначально перед кластерами ставились две задачи: достичь большой вычислительной мощности и обеспечить повышенную надёжность ВС. Архитектура кластерных систем во многом похожа на архитектуру MPP-систем: тот же принцип распределённой памяти, использование в качестве вычислительных узлов законченных вычислительных машин, большой потенциал для масштабирования. В качестве узла кластера может выступать как однопроцессорная ВМ, так и ВС типа SMP. Как правило, это не специализированные устройства, как в MPP, а серийно выпускаемые вычислительные машины и системы. Кластерные системы с одинаковыми узлами называют гомогенными, а с разнотипными узлами — гетерогенными. Использование машин массового производства существенно снижает стоимость ВС, а возможность варьирования различных по типу узлов позволяет получить необходимую производительность за приемлемую цену. Узлы в кластерной системе объединены высокоскоростной сетью. Неотъемлемая часть кластера — специализированное программное обеспечение, организующее бесперебойную работу при отказе одного или нескольких узлов. Кластеры обеспечивают высокий уровень доступности — в них отсутствуют единая операционная система и совместно используемая память, то есть нет проблемы когерентности кэшей. Кластеры хорошо масштабируются путём добавления узлов, что позволяет достичь высочайших показателей производительности. Примеры кластерных систем: Roadrunner BladeCenter QS22 (IBM) с пиковой производительностью 1376 TFLOPS, состоящая из 122\,400 узлов; многочисленные кластеры на базе Ethernet и SCI.
\textit{(Источник: Орлов, Цилькер, гл. 14, с. 569--572.)}

\paragraph{Кластеры больших SMP-систем (Constellations).}
Огромный потенциал масштабирования, свойственный кластерной архитектуре, делает её очень перспективным направлением в области создания высокопроизводительных вычислительных систем. Масштабирование возможно как за счёт увеличения числа узлов, так и путём применения в качестве узлов не одиночных ВМ, а хорошо масштабируемых вычислительных систем, обычно SMP-типа. Это направление получило настолько широкое развитие, что было выделено в отдельную группу MIMD-устройств — Constellations (термин, обусловленный названием одной из первых ВС данного типа — Sun «Constellation», можно перевести как «созвездие»). Constellation-система — это кластер, узлами которого служат SMP-системы. В качестве отличительного признака выступает число процессорных элементов в узле: в настоящее время условием причисления ВС к Constellation-системам служит число ПЭ в узле — оно должно быть больше (равно) 16; системы, не отвечающие данному условию, считаются классическими кластерными системами. Перспективность Constellation-систем обусловлена удачным сочетанием преимуществ распределённой памяти (возможности наращивания количества узлов) и разделяемой памяти (эффективного доступа множества ПЭ к памяти). Формально структура Constellation-системы полностью соответствует кластерной архитектуре, однако явная направленность на высокую производительность может отражаться в некоторых конструктивных решениях. Пример: Tera 10 фирмы Bull с пиковой производительностью 58,8 TFLOPS, кластер из 544 узлов, каждый узел — SMP-система из 8 двухъядерных процессоров Itanium 2.
\textit{(Источник: Орлов, Цилькер, гл. 14, с. 572--573.)}

\paragraph{Многомашинные комплексы.}
Многомашинный комплекс представляет собой совокупность отдельных вычислительных машин, объединённых коммуникационной сетью для совместного решения задач. Каждая машина в комплексе сохраняет самостоятельность и может работать под управлением собственной операционной системы. Взаимодействие между машинами осуществляется путём обмена сообщениями. Многомашинные комплексы могут строиться на базе серийно выпускаемых ВМ и использоваться для повышения производительности, надёжности или для решения задач, требующих распределённой обработки. К многомашинным комплексам относятся, в частности, кластеры и Constellation-системы. Основное различие между многопроцессорными и многомашинными системами заключается в степени связи: многопроцессорные системы с разделяемой памятью являются сильно связанными, а многомашинные комплексы — слабо связанными. В многомашинных комплексах каждый узел имеет собственную память и не может напрямую обращаться к памяти другого узла; для обмена данными используется передача сообщений по сети. Это упрощает масштабирование и повышает отказоустойчивость, но увеличивает накладные расходы на коммуникации.
\textit{(Источник: Орлов, Цилькер, гл. 11, с. 466--467; гл. 14, с. 557--578.)}

\subsubsection{Сравнительный анализ архитектур}

\paragraph{Сводная таблица.}
В таблице \ref{tab:comparison} приведено сравнение основных архитектур многопроцессорных и многомашинных комплексов по ключевым параметрам.

\begin{table}[h]
\centering
\caption{Сравнение архитектур многопроцессорных и многомашинных комплексов}
\label{tab:comparison}
\begin{tabular}{p{3cm}p{3cm}p{3cm}p{3cm}}
\toprule
\textbf{Архитектура} & \textbf{Память} & \textbf{Масштабируемость} & \textbf{Особенности} \\
\midrule
SMP & Разделяемая (UMA) & До 32--64 процессоров & Простота, но ограниченная масштабируемость; когерентность кэшей аппаратная \\
NUMA & Разделяемая (NUMA) & До 1024 процессоров & Неоднородный доступ; требует протоколов справочника \\
MPP & Распределённая (NORMA) & Тысячи узлов & Высокая производительность; хост-компьютер — узкое место \\
Кластеры & Распределённая (NORMA) & Тысячи узлов & Высокая доступность; гомогенные/гетерогенные; стандартные компоненты \\
Constellations & Распределённая (NORMA) & Сотни узлов & Кластер из SMP-узлов; сочетание преимуществ \\
\bottomrule
\end{tabular}
\end{table}

\paragraph{Области применения.}
SMP-системы применяются для задач, требующих интенсивных вычислений с частым обменом данными между потоками, например в научных расчётах, базах данных, серверах приложений. NUMA-системы используются для крупномасштабных задач, где требуется большой объём памяти и высокая производительность, например в суперкомпьютерах. MPP-системы ориентированы на задачи с массовым параллелизмом, такие как моделирование климата, физика высоких энергий, геномика. Кластеры находят широкое применение в качестве недорогих суперкомпьютеров, веб-серверов, систем обработки данных. Constellation-системы используются для построения высокопроизводительных вычислительных комплексов на базе стандартных SMP-узлов.
\textit{(Источник: Орлов, Цилькер, гл. 14, с. 557--578.)}

\subsubsection{Краткое резюме подраздела}
Многопроцессорные и многомашинные комплексы классифицируются по способу организации памяти на системы с разделяемой памятью (сильно связанные) и системы с распределённой памятью (слабо связанные). К системам с разделяемой памятью относятся SMP (UMA) и NUMA (ccNUMA, nccNUMA). К системам с распределённой памятью — MPP, кластеры и Constellation-системы. Каждая архитектура имеет свои достоинства и недостатки, определяющие область её применения. SMP-системы просты, но ограниченно масштабируемы; NUMA обеспечивает масштабируемость до тысяч процессоров, но требует аппаратной поддержки когерентности; MPP и кластеры строятся из независимых узлов и обеспечивают высокую производительность и отказоустойчивость; Constellation-системы сочетают преимущества кластеров и SMP. Выбор конкретной архитектуры определяется требованиями к производительности, масштабируемости, стоимости и надёжности.
\textit{(Обобщение по: Орлов, Цилькер, гл. 11, с. 466--492; гл. 14, с. 557--578.)}

\subsection{Вычислительные кластеры}

\noindent\textit{Основные источники: Орлов С.А., Цилькер Б.Я. Организация ЭВМ и систем: Учебник для вузов. 2-е изд. СПб.: Питер, 2011. Гл. 14, с. 569--573; гл. 11, с. 466--467.}

\subsubsection{Определение и назначение кластеров}

\paragraph{Понятие кластера.}
Кластером называют группу взаимно соединённых вычислительных систем (узлов), работающих совместно и составляющих единый вычислительный ресурс, создавая иллюзию наличия единственной ВМ. Изначально перед кластерами ставились две задачи: достичь большой вычислительной мощности и обеспечить повышенную надёжность ВС. Пионером в области кластерных архитектур считается корпорация DEC, разработавшая первый коммерческий кластер в начале 80-х годов прошлого века.
\textit{(Источник: Орлов, Цилькер, гл. 14, с. 569.)}

\paragraph{Место кластеров в классификации MIMD-систем.}
Архитектура кластерных систем во многом похожа на архитектуру MPP-систем: тот же принцип распределённой памяти, использование в качестве вычислительных узлов законченных вычислительных машин, большой потенциал для масштабирования системы и целый ряд других особенностей. В первом приближении кластерную технологию можно рассматривать как развитие идей массовых параллельных вычислений. С другой стороны, многие черты кластерной архитектуры дают основание считать её самостоятельным направлением в области MIMD-систем. Кластеры относятся к классу систем с распределённой памятью (слабо связанных, NORMA), в которых каждый узел имеет доступ только к своей локальной памяти, а взаимодействие осуществляется путём обмена сообщениями по сети.
\textit{(Источник: Орлов, Цилькер, гл. 14, с. 569.)}

\subsubsection{Архитектура кластера}

\paragraph{Узлы кластера.}
В качестве узла кластера может выступать как однопроцессорная ВМ, так и ВС типа SMP (напомним, что логически SMP-система представляется как единственная ВМ). Как правило, это не специализированные устройства, приспособленные под использование в вычислительной системе, как в MPP, а серийно выпускаемые вычислительные машины и системы. Ещё одна особенность кластерной архитектуры состоит в том, что в единую систему объединяются узлы разного типа, от персональных компьютеров до мощных ВС. Кластерные системы с одинаковыми узлами называют \emph{гомогенными} кластерами, а с разнотипными узлами — \emph{гетерогенными} кластерами. Использование машин массового производства существенно снижает стоимость ВС, а возможность варьирования различных по типу узлов позволяет получить необходимую производительность за приемлемую цену. Важно и то, что узлы могут функционировать самостоятельно и отдельно от кластера; для этого каждый узел работает под управлением своей операционной системы. Чаще всего используются стандартные ОС: Linux, FreeBSD, Solaris и версии Windows, продолжающие направление Windows NT.
\textit{(Источник: Орлов, Цилькер, гл. 14, с. 569--570.)}

\paragraph{Сеть и протоколы.}
Узлы в кластерной системе объединены высокоскоростной сетью. Решения могут быть простыми, основанными на аппаратуре Ethernet, или сложными, с высокоскоростными сетями пропускной способности в сотни мегабайтов в секунду. К последней категории относятся сети SCI компании Scali Computer (около 100 Мбит/с) и Myrinet (около 120 Мбит/с). В принципе, за основу кластерной системы может быть взята стандартная локальная сеть (или сеть большего масштаба) с сохранением принятых в ней протоколов. Аппаратурные изменения могут не потребоваться или, в худшем случае, сводятся к замене коммуникационного оборудования на более производительное. При соединении машин в кластер почти всегда поддерживаются прямые межмашинные связи. Вычислительные машины (системы) в кластере взаимодействуют в соответствии с одним из двух транспортных протоколов: TCP (Transmission Control Protocol) или UDP (User Datagram Protocol). В последнее время применяют специальные протоколы, которые работают намного лучше, например Virtual Interface Architecture (VIA).
\textit{(Источник: Орлов, Цилькер, гл. 14, с. 570.)}

При обмене информацией используются два программных метода: передачи сообщений и распределённой совместно используемой памяти. Первый опирается на явную передачу информационных сообщений между узлами кластера. В альтернативном варианте также происходит пересылка сообщений, но движение данных между узлами кластера скрыто от программиста. Подключение узлов к сети осуществляется посредством сетевых адаптеров. Учитывая роль коммуникаций, для связи ядра вычислительного узла с сетью используется выделенная шина ввода/вывода; к этой шине также подключаются локальные магнитные диски (наличие таких дисков типично для кластерных систем, но не характерно для MPP-систем). Элементы вычислительного ядра объединяются посредством локальной системной шины; связь между этой шиной и шиной ввода/вывода обеспечивает мост.
\textit{(Источник: Орлов, Цилькер, гл. 14, с. 570.)}

\subsubsection{Программное обеспечение кластера}

\paragraph{Кластерное ПО.}
Неотъемлемая часть кластера — специализированное программное обеспечение (ПО), организующее бесперебойную работу при отказе одного или нескольких узлов. Такое ПО должно быть установлено на каждый узел кластера. Оно реализует механизм передачи сообщений над стандартными сетевыми протоколами и может рассматриваться как часть операционной системы. Именно благодаря специализированному ПО группа ВМ, объединённых сетью, превращается в кластерную вычислительную систему. Кластерное ПО перераспределяет вычислительную нагрузку при отказе одного или нескольких узлов кластера, а также восстанавливает вычисления при сбое в узле.
\textit{(Источник: Орлов, Цилькер, гл. 14, с. 570--571.)}

\paragraph{Механизм контроля работоспособности.}
ПО каждого узла постоянно контролирует работоспособность всех остальных узлов на основе периодической рассылки каждым узлом сигнала, известного как keepalive («пока жив») или heartbeat («сердцебиение»). Если сигнал от некоторого узла не поступает, узел считается вышедшим из строя; ему не даётся возможность выполнять ввод/вывод, его диски и другие ресурсы (включая сетевые адреса) переназначаются другим узлам, а выполнявшиеся им программы перезапускаются в других узлах. При наличии в кластере совместно используемых дисков кластерное ПО поддерживает единую файловую систему.
\textit{(Источник: Орлов, Цилькер, гл. 14, с. 571.)}

\paragraph{Подходы к построению кластеров.}
При создании кластерных систем используется один из двух подходов. Первый подход применяется для построения небольших кластерных систем, например на базе небольших локальных сетей организаций или их подразделений. В таком кластере каждая ВМ продолжает работать как самостоятельная единица, одновременно выполняя функции узла кластерной системы. Второй подход ориентирован на использование кластерной системы в роли мощного вычислительного ресурса. Узлами кластера служат только системные блоки вычислительных машин, компактно размещаемые в специальных стойках. Управление системой и запуск задач осуществляет полнофункциональный хост-компьютер; он же поддерживает дисковую подсистему кластера и разнообразное периферийное оборудование. Отсутствие у узлов собственной периферии существенно удешевляет систему.
\textit{(Источник: Орлов, Цилькер, гл. 14, с. 571.)}

\subsubsection{Классификация кластеров}

\paragraph{По назначению.}
Кластеры можно классифицировать по их назначению. \emph{Высокопроизводительные кластеры} (High Performance Computing, HPC) ориентированы на максимальную вычислительную мощность и используются для решения сложных научных и инженерных задач. \emph{Кластеры высокой доступности} (High Availability, HA) предназначены для обеспечения бесперебойной работы сервисов; они обеспечивают резервирование и автоматическое переключение на исправный узел при отказе. \emph{Кластеры балансировки нагрузки} распределяют запросы между узлами для повышения производительности и надёжности веб-сервисов. Возможны также комбинированные варианты.
\textit{(Источник: Орлов, Цилькер, гл. 14, с. 569--571.)}

\paragraph{По типу узлов.}
По типу узлов различают \emph{гомогенные} кластеры, в которых все узлы одинаковы, и \emph{гетерогенные} кластеры, в которых узлы различаются по типу и производительности. Гомогенные кластеры проще в настройке и управлении, гетерогенные позволяют более гибко использовать имеющиеся ресурсы.
\textit{(Источник: Орлов, Цилькер, гл. 14, с. 569.)}

\subsubsection{Преимущества и недостатки кластеризации}

\paragraph{Преимущества.}
В работе [59] перечисляются четыре преимущества, достигаемые с помощью кластеризации:
\begin{itemize}
\item \emph{абсолютная масштабируемость}: возможно создание больших кластеров, превосходящих по вычислительной мощности даже самые производительные одиночные ВМ; кластер в состоянии содержать десятки узлов, каждый из которых представляет собой мультипроцессор;
\item \emph{наращиваемая масштабируемость}: кластер строится так, что его можно наращивать, добавляя новые узлы небольшими порциями; пользователь может начать с умеренной системы, расширяя её по мере необходимости;
\item \emph{высокий коэффициент готовности}: поскольку каждый узел кластера — самостоятельная ВМ или ВС, отказ одного из узлов не приводит к потере работоспособности кластера; во многих системах отказоустойчивость автоматически поддерживается программным обеспечением;
\item \emph{превосходное соотношение цена/производительность}: кластер любой производительности можно создать, соединяя стандартные «строительные блоки», при этом его стоимость будет ниже, чем у одиночной ВМ с эквивалентной вычислительной мощностью.
\end{itemize}
\textit{(Источник: Орлов, Цилькер, гл. 14, с. 571.)}

\paragraph{Недостатки.}
Основной недостаток кластерных систем — взаимодействие между узлами кластера занимает гораздо больше времени, чем в других типах ВС. Это связано с тем, что для обмена данными используется сеть, а не общая память; передача сообщений требует дополнительных накладных расходов. Кроме того, кластеры, как правило, не обеспечивают автоматического распараллеливания программ; программист должен явно использовать средства обмена сообщениями (например, MPI) или распределённую память. Тем не менее кластеры хорошо масштабируются путём добавления узлов, что позволяет достичь высочайших показателей производительности; благодаря этой особенности архитектуры кластеры с сотнями и тысячами узлов положительно зарекомендовали себя на практике.
\textit{(Источник: Орлов, Цилькер, гл. 14, с. 571--572.)}

\subsubsection{Примеры кластерных систем}

\paragraph{Roadrunner.}
До недавнего времени именно кластерная ВС занимала первую позицию в списке самых производительных систем. Речь идёт о кластерной системе Roadrunner BladeCenter QS22, созданной компанией IBM. Теоретическая пиковая производительность системы составляет 1376 TFLOPS. Система состоит из 122\,400 узлов на базе процессоров Opteron и PowerXCell. Roadrunner использовался для решения задач ядерной безопасности и других научных расчётов.
\textit{(Источник: Орлов, Цилькер, гл. 14, с. 572.)}

\paragraph{Другие примеры.}
Кроме Roadrunner, в список Top500 входят и другие кластерные системы: Tianhe-1 (Китай), Ranger (США), Red Sky (США). Кластерные системы строятся на базе стандартных компонентов и часто используются в научных и образовательных учреждениях. Например, кластер из 544 SMP-узлов Tera 10 фирмы Bull с пиковой производительностью 58,8 TFLOPS относится к Constellation-системам (см. следующий подраздел). Многие кластеры используют сети Ethernet, SCI, Myrinet и InfiniBand.
\textit{(Источник: Орлов, Цилькер, гл. 14, с. 572.)}

\subsubsection{Сравнение с MPP}

\paragraph{Сходство и различие.}
Кластеры и MPP-системы имеют много общего: распределённая память, использование законченных вычислительных машин в качестве узлов, высокий потенциал масштабирования. Однако есть и существенные различия. В MPP-системах узлы, как правило, специализированы, а система содержит выделенный хост-компьютер; в кластерах используются серийные машины, и каждый узел может работать под управлением собственной ОС. Кластеры обеспечивают более высокий коэффициент готовности за счёт программного обеспечения, в то время как в MPP отказ хост-компьютера может привести к отказу всей системы. Кластеры обычно дешевле MPP и проще масштабируются, но уступают MPP в максимальной производительности и плотности упаковки.
\textit{(Источник: Орлов, Цилькер, гл. 14, с. 569--572.)}

\subsubsection{Кластеры больших SMP-систем (Constellations)}

\paragraph{Определение и особенности.}
Огромный потенциал масштабирования, свойственный кластерной архитектуре, делает её очень перспективным направлением в области создания высокопроизводительных вычислительных систем. Масштабирование возможно как за счёт увеличения числа узлов, так и путём применения в качестве узлов не одиночных ВМ, а хорошо масштабируемых вычислительных систем, обычно SMP-типа. Это направление получило настолько широкое развитие, что было выделено в отдельную группу MIMD-устройств — Constellations (термин, обусловленный названием одной из первых ВС данного типа — Sun «Constellation», можно перевести как «созвездие»). Constellation-система — это кластер, узлами которого служат SMP-системы. В качестве отличительного признака выступает число процессорных элементов в узле: в настоящее время условием причисления ВС к Constellation-системам служит число ПЭ в узле — оно должно быть больше (равно) 16; системы, не отвечающие данному условию, считаются классическими кластерными системами. Перспективность Constellation-систем обусловлена удачным сочетанием преимуществ распределённой памяти (возможности наращивания количества узлов) и разделяемой памяти (эффективного доступа множества ПЭ к памяти). Формально структура Constellation-системы полностью соответствует кластерной архитектуре, однако явная направленность на высокую производительность может отражаться в некоторых конструктивных решениях.
\textit{(Источник: Орлов, Цилькер, гл. 14, с. 572--573.)}

\paragraph{Пример.}
В качестве примера Constellation-системы можно привести систему Tera 10 фирмы Bull с теоретической пиковой производительностью 58,8 TFLOPS. Система представляет собой кластер из 544 узлов, в котором каждый узел — это SMP-система, образованная 8 двухъядерными процессорами Itanium 2.
\textit{(Источник: Орлов, Цилькер, гл. 14, с. 573.)}

\subsubsection{Краткое резюме подраздела}
Кластер — это группа взаимно соединённых вычислительных систем (узлов), работающих совместно как единый ресурс. Кластеры относятся к классу MIMD-систем с распределённой памятью (слабо связанных). Узлами кластера служат серийно выпускаемые ВМ или SMP-системы; кластеры могут быть гомогенными и гетерогенными. Сеть кластера строится на основе Ethernet, SCI, Myrinet и других технологий; для обмена сообщениями используются протоколы TCP, UDP, VIA. Кластерное ПО обеспечивает контроль работоспособности узлов (keepalive/heartbeat), перераспределение нагрузки и восстановление вычислений при отказах. Кластеры обладают абсолютной и наращиваемой масштабируемостью, высоким коэффициентом готовности и хорошим соотношением цена/производительность. Основной недостаток — большие коммуникационные задержки. Кластеры хорошо масштабируются и могут достигать высочайшей производительности; пример — Roadrunner (IBM) с пиковой производительностью 1376 TFLOPS. Развитием кластерной архитектуры являются Constellation-системы, узлами которых служат SMP-системы с числом ПЭ не менее 16.
\textit{(Обобщение по: Орлов, Цилькер, гл. 14, с. 569--573.)}

\subsection{Проблемно-ориентированные параллельные структуры}

\noindent\textit{Основные источники: Орлов С.А., Цилькер Б.Я. Организация ЭВМ и систем: Учебник для вузов. 2-е изд. СПб.: Питер, 2011. Гл. 13, с. 526--556; заключение, с. 598.}

\subsubsection{Общая характеристика}

\paragraph{Назначение и место в классификации.}
Проблемно-ориентированные параллельные структуры представляют собой специализированные вычислительные системы, архитектура которых оптимизирована под определённый класс задач. В отличие от универсальных MIMD-систем, они ориентированы на эффективное выполнение операций над массивами данных, матричных вычислений, цифровой обработки сигналов и других задач с регулярной структурой. К проблемно-ориентированным параллельным структурам относятся матричные вычислительные системы, систолические структуры и нейросети. Согласно классификации Флинна, матричные и систолические системы обычно относят к классу SIMD (одиночный поток команд — множественный поток данных), хотя систолические структуры могут сочетать параллелизм с конвейерной обработкой. Нейросети в зависимости от реализации могут тяготеть к SIMD или MIMD.
\textit{(Источник: Орлов, Цилькер, гл. 13, с. 526--556; заключение, с. 598.)}

\subsubsection{Матричные вычислительные системы}

\paragraph{Общая структура.}
Матричная SIMD-система содержит массив процессорных элементов (МПЭ) и контроллер массива процессорных элементов (КМП). КМП выполняет последовательную часть программы, декодирует команды и транслирует в массив управляющие сигналы, команды и данные. Каждый процессорный элемент имеет локальную память и обрабатывает свои данные по общей команде. Для связи с пользователем используется фронтальная ВМ (front-end computer). Структура матричной системы показана на рис. 13.12 (см. учебник).
\textit{(Источник: Орлов, Цилькер, гл. 13, с. 535--537.)}

\paragraph{Контроллер массива процессорных элементов.}
КМП выполняет последовательный программный код, реализует команды ветвления программы, транслирует команды и сигналы управления в процессорные элементы. Одна из возможных реализаций КМП, принятая в устройстве управления матричной вычислительной системы PASM, показана на рис. 13.13 (см. учебник). При загрузке из фронтальной ВМ программа через интерфейс ввода/вывода заносится в оперативное запоминающее устройство контроллера (ОЗУ КМП). Команды для процессорных элементов и глобальная маска, формируемая на этапе компиляции, также через интерфейс ввода/вывода загружаются в ОЗУ команд и глобальной маски (ОЗУ КГМ). Затем контроллер начинает выполнять программу, извлекая либо одну скалярную команду из ОЗУ КМП, либо множественные команды из ОЗУ КГМ. Скалярные команды выполняются центральным процессором (ЦП) контроллера; команды, оперирующие параллельными переменными, преобразуются в блоке выборки команд в нанокоманды. Нанокоманды совместно с маской пересылаются через шину команд для ПЭ на исполнение в массив процессорных элементов. Например, команда сложения 32-разрядных слов в КМП системы MPP преобразуется в 32 нанокоманды одноразрядного сложения, которые каждым ПЭ обрабатываются последовательно. Для обеспечения ветвления программы статусная информация, хранящаяся в ПЭ, собирается в единое слово и передаётся в КМП по шине результата.
\textit{(Источник: Орлов, Цилькер, гл. 13, с. 537--538.)}

\paragraph{Структура процессорного элемента.}
В большинстве матричных SIMD-систем в качестве процессорных элементов применяются простые RISC-процессоры с локальной памятью ограниченной ёмкости. Например, каждый ПЭ системы MasPar MP-1 состоит из четырёхразрядного процессора с памятью ёмкостью 64 Кбайт; в системе MPP используются одноразрядные процессоры с памятью 1 Кбит; в CM-2 процессорный элемент представляет собой одноразрядный процессор с 64 Кбит локальной памяти. Неотъемлемыми компонентами ПЭ (рис. 13.15, см. учебник) являются: арифметико-логическое устройство (АЛУ); регистры данных; сетевой интерфейс (СИ), который может включать регистры пересылки данных; номер процессора; регистр флага разрешения маскирования \(F\); локальная память. Процессорные элементы, управляемые командами из КМП, могут выбирать данные из своей локальной памяти и регистров, обрабатывать их в АЛУ и сохранять результаты. Обмен данными между ПЭ производится по коммуникационной сети; для передачи результатов в КМП служит шина результата.
\textit{(Источник: Орлов, Цилькер, гл. 13, с. 539--541.)}

\paragraph{Маскирование.}
В процессе вычислений в ряде операций должны участвовать только определённые ПЭ, в то время как остальные остаются бездействующими. Разрешение и запрет работы ПЭ могут исходить от контроллера массива (глобальное маскирование) и реализуются с помощью схем маскирования. В этом случае решение о необходимости маскирования принимается на этапе компиляции кода. Решение о маскировании может также приниматься во время выполнения программы (маскирование, определяемое данными), при этом опираются на хранящийся в ПЭ флаг разрешения маскирования \(F\). При использовании схемы глобального маскирования КМП вместе с командами посылает во все ПЭ глобальную маску; каждый ПЭ декодирует её и выясняет, должен ли он выполнять данную команду. Различают маску в виде \(N\)-разрядного вектора (каждый бит отражает состояние одного ПЭ) и маску с адресом ПЭ (позиции соответствуют разрядам двоичного адреса ПЭ, допускаются значения 0, 1, X). Глобальные и локальные схемы маскирования могут комбинироваться: активность ПЭ определяется как флагом \(F\), так и глобальной маской.
\textit{(Источник: Орлов, Цилькер, гл. 13, с. 541--542.)}

\paragraph{Сети взаимосвязей.}
Эффективность матричной системы во многом определяется топологией сети взаимосвязей процессорных элементов. Наиболее распространены решётчатые и гиперкубические топологии. Так, в ILLIAC IV, MPP и CM-2 каждый ПЭ соединён с четырьмя соседними; в MP-1 и MP-2 каждый ПЭ связан с восемью смежными ПЭ. В ряде систем (MP-1, MP-2, GF11) реализуются многоступенчатые динамические сети соединений. При синхронной работе ПЭ обмен информацией должен производиться по согласованной схеме; для предотвращения потери данных при одновременной передаче от нескольких источников к одному приёмнику используются специальные механизмы (например, объединение сообщений в CM-2, возможность записи в разные ячейки локальной памяти в MP-1).
\textit{(Источник: Орлов, Цилькер, гл. 13, с. 542--543.)}

\paragraph{Достоинства и недостатки.}
Матричные процессоры обеспечивают высокий уровень параллелизма при выполнении однотипных операций над большими массивами данных. Они эффективны для задач, хорошо распараллеливаемых по данным (обработка изображений, решение сеточных задач, матричные вычисления). Основные недостатки: сложность программирования, необходимость согласования данных, ограниченная применимость для задач с нерегулярной структурой, а также высокая стоимость при большом числе ПЭ.
\textit{(Источник: Орлов, Цилькер, гл. 13, с. 535--543.)}

\subsubsection{Систолические структуры}

\paragraph{Идея систолической обработки.}
В фон-неймановских машинах данные, считанные из памяти, однократно обрабатываются в процессорном элементе, после чего снова возвращаются в память. Авторы идеи систолической матрицы Кунг и Лейзерсон предложили организовать вычисления так, чтобы данные на своём пути от считывания из памяти до возвращения обратно пропускались через как можно большее число ПЭ. Если проводить аналогию между структурами ВС и живого организма, то памяти можно отвести роль сердца, множеству ПЭ — роль тканей, а поток данных следует рассматривать как циркулирующую кровь. Отсюда и происходит название «систолическая матрица» (систола — сокращение предсердий и желудочков сердца, при котором кровь нагнетается в артерии). Систолические структуры эффективны при выполнении матричных вычислений, обработке сигналов, сортировке данных.
\textit{(Источник: Орлов, Цилькер, гл. 13, с. 547.)}

\paragraph{Процессорный элемент и принцип работы.}
Основу систолической структуры образует однородная вычислительная среда из процессорных элементов. Типичный ПЭ реализует операцию «умножение с накоплением» (IPS-элемент). На вход ПЭ подаются операнды \(a_{вх}\), \(b_{вх}\); на выходе формируются \(a_{вых}\), \(b_{вых}\) и частичная сумма \(c_{вых}\). На \(n\)-м шаге работы систолической системы ПЭ выполняет операцию
\[
c^{(n)}_{вых}=c^{(n-1)}_{вх}+a^{(n-1)}_{вх}\times b^{(n-1)}_{вх},
\]
при этом операнды на входе и выходе ПЭ одинаковы: \(a^{(n)}_{вых}=a^{(n-1)}_{вх}\), \(b^{(n)}_{вых}=b^{(n-1)}_{вх}\). Частичная сумма поступает на вход ПЭ либо с данного процессорного элемента, либо с соседнего ПЭ матрицы. Вычислительный процесс в систолических структурах представляет собой непрерывную и регулярную передачу данных от одного ПЭ к другому без запоминания промежуточных результатов; каждый элемент входных данных выбирается из памяти однократно и используется столько раз, сколько необходимо по алгоритму; ввод данных осуществляется в крайние ПЭ матрицы. Функциональная схема IPS-элемента показана на рис. 13.23 (см. учебник).
\textit{(Источник: Орлов, Цилькер, гл. 13, с. 548--549.)}

\paragraph{Классификация систолических структур.}
Систолические структуры классифицируются по нескольким признакам:
\begin{itemize}
\item по степени гибкости: специализированные (ориентированы на определённый алгоритм), алгоритмически ориентированные (на класс алгоритмов), программируемые (с возможностью программирования ПЭ и конфигурации связей);
\item по разрядности ПЭ: одноразрядные и многоразрядные;
\item по характеру локально-пространственных связей: одномерные (линейные), двухмерные, трёхмерные.
\end{itemize}
Наиболее распространены линейные матрицы (ULA, BLA, TLA). ULA (Unidirectional Linear Array) — однонаправленный линейный массив, где поток данных перемещается в одном направлении. BLA (Bidirectional Linear Array) — двунаправленный линейный массив, в котором два потока данных движутся навстречу друг другу. TLA (Three-path communication Linear Array) — линейный процессорный массив с тремя коммуникационными трактами, в котором по разным направлениям перемещаются три потока данных. Выбор конфигурации зависит от вида обрабатываемой информации: одномерные схемы применяются при обработке векторов, двухмерные — матриц, трёхмерные — коллекций иного типа.
\textit{(Источник: Орлов, Цилькер, гл. 13, с. 549--552.)}

\paragraph{Пример вычислений.}
Организацию вычислительного процесса в систолических массивах удобно пояснить на примере умножения матрицы \(A=\{a_{ij}\}\) на вектор \(X=\{x_1,x_2,\ldots,x_n\}\). Элементы вектора произведения \(Y=\{y_1,y_2,\ldots,y_n\}\) могут быть получены периодически повторяющимися операциями
\[
y_i(1)=0;\qquad y_i(k+1)=y_i(k)+a_i(k)\times x_i(k);\qquad y_i=y_i(n+1),
\]
где \(k\) — номер шага вычислений. Компоненты \(y_i\) вектора \(Y\), имеющие в начальный момент нулевое значение, поступают на вход массива и продвигаются через ПЭ справа налево, в то время как компоненты вектора \(X\) движутся слева направо. Элементы матрицы \(A=\{a_{ij}\}\) в порядке, определяемом схемой, вводятся в ПЭ сверху вниз. Промежуточные результаты \(y_i(k)\) накапливаются по мере продвижения от одного ПЭ к другому. В каждый момент времени активны только \(\frac{p+q-1}{2}\) процессорных элементов (где \(p+q-1\) — ширина полосы ненулевых элементов матрицы), поэтому каждый выходной результат формируется за два такта. Для вычисления всех \(n\) элементов выходного вектора \(Y\) необходимо \(2n+p+q-1\) тактов.
\textit{(Источник: Орлов, Цилькер, гл. 13, с. 553--555.)}

\paragraph{Достоинства и недостатки.}
Систолические структуры обладают следующими особенностями:
\begin{itemize}
\item вычислительный процесс представляет собой непрерывную и регулярную передачу данных от одного ПЭ к другому без запоминания промежуточных результатов;
\item каждый элемент входных данных выбирается из памяти однократно и используется многократно;
\item образующие структуру ПЭ однотипны и могут быть менее универсальными, чем процессоры обычных многопроцессорных систем;
\item потоки данных и управляющих сигналов обладают регулярностью, что позволяет объединять ПЭ локальными связями минимальной длины;
\item алгоритмы функционирования позволяют совместить параллелизм с конвейерной обработкой данных;
\item производительность матрицы можно улучшить за счёт добавления определённого числа ПЭ, причём коэффициент повышения производительности при этом линеен.
\end{itemize}
К недостаткам можно отнести сложность программирования и ограниченную применимость для задач с нерегулярной структурой.
\textit{(Источник: Орлов, Цилькер, гл. 13, с. 549.)}

\subsubsection{Нейросети}

\noindent\textit{Основные источники: Хливненко Л.В., Пятакович Ф.А. Практика нейросетевого моделирования: учебное пособие. СПб.: Лань, 2023. 200 с.; Хайкин С. Нейронные сети: полный курс. 2-е изд. М.: Диалектика, 2020. 1103 с.}

\paragraph{Понятие искусственной нейронной сети.}
Искусственная нейронная сеть (ИНС) представляет собой параллельную вычислительную структуру, состоящую из множества простых обрабатывающих элементов — искусственных нейронов, соединённых между собой взвешенными связями (синапсами). Каждый нейрон получает входные сигналы от других нейронов, суммирует их с учётом весов связей и преобразует результат с помощью нелинейной функции активации. Нейронная сеть может быть реализована аппаратно (нейропроцессоры, нейрочипы) или программно (на универсальных процессорах, графических ускорителях, специализированных ПЛИС). Основные свойства нейросетей — способность к обучению, обобщению, параллельной обработке информации и устойчивость к ошибкам. Нейронные сети, или, точнее, искусственные нейронные сети, представляют собой технологию, уходящую корнями во множество дисциплин: нейрофизиологию, математику, статистику, физику, компьютерные науки и технику. Они находят применение в таких разнородных областях, как моделирование, анализ временных рядов, распознавание образов, обработка сигналов и управление, благодаря одному важному свойству — способности обучаться на основе данных при участии учителя или без его вмешательства.
\textit{(Источник: Хайкин, 2020, аннотация; Хливненко, Пятакович, 2023, с. 3--5.)}

\paragraph{Искусственный нейрон.}
Основным элементом нейронной сети является искусственный нейрон. Он состоит из трёх основных частей: набора синапсов (входных связей), каждый из которых характеризуется своим весом; сумматора, складывающего взвешенные входные сигналы; и функции активации, которая ограничивает амплитуду выходного сигнала. Математически нейрон описывается выражением
\[
y = \varphi\left(\sum_{i=1}^{n} w_i x_i + b\right),
\]
где \(x_i\) — входные сигналы, \(w_i\) — веса связей, \(b\) — порог (смещение), \(\varphi\) — функция активации. В качестве функций активации обычно используются пороговая, сигмоидальная, гиперболический тангенс, ReLU и другие. Выбор функции активации существенно влияет на свойства сети и процесс её обучения.
\textit{(Источник: Хайкин, 2020, гл. 1; Хливненко, Пятакович, 2023, с. 5--8.)}

\paragraph{Архитектуры нейронных сетей.}
Различают несколько основных архитектур нейронных сетей:
\begin{itemize}
\item \emph{сети прямого распространения} (перцептроны): сигнал распространяется от входного слоя к выходному без обратных связей; простейший пример — многослойный перцептрон (MLP), содержащий один или несколько скрытых слоёв;
\item \emph{рекуррентные сети} (RNN): содержат обратные связи, что позволяет хранить информацию о предыдущих состояниях и обрабатывать последовательности; примеры — сети Хопфилда, LSTM (Long Short-Term Memory);
\item \emph{свёрточные сети} (CNN): используют операции свёртки для обработки изображений и других данных с регулярной структурой; состоят из чередующихся свёрточных и субдискретизирующих (pooling) слоёв, а также полносвязных слоёв на выходе;
\item \emph{сети с радиальными базисными функциями} (RBF), \emph{самоорганизующиеся карты Кохонена}, \emph{автоэнкодеры} и другие.
\end{itemize}
Каждая архитектура ориентирована на определённый класс задач: распознавание образов, обработка естественного языка, прогнозирование, управление, сжатие данных.
\textit{(Источник: Хливненко, Пятакович, 2023, с. 8--15; Хайкин, 2020, гл. 4--8.)}

\paragraph{Обучение нейронных сетей.}
Обучение нейросети обычно сводится к настройке весов связей по обучающей выборке. Для этого используются алгоритмы обратного распространения ошибки, стохастического градиентного спуска и другие. В процессе обучения сеть адаптируется к входным данным и вырабатывает способность к обобщению на новые примеры. Различают:
\begin{itemize}
\item \emph{обучение с учителем} (supervised learning): для каждого входного примера задан желаемый выход (целевое значение); сеть подстраивает веса так, чтобы минимизировать ошибку между фактическим и желаемым выходом;
\item \emph{обучение без учителя} (unsupervised learning): целевые значения не заданы; сеть выявляет скрытые структуры в данных, выполняет кластеризацию или самоорганизацию;
\item \emph{обучение с подкреплением} (reinforcement learning): сеть обучается на основе сигнала вознаграждения, получаемого за действия в некоторой среде.
\end{itemize}
Для обучения глубоких сетей применяются методы регуляризации (dropout, batch normalization), а также специальные архитектуры и алгоритмы оптимизации.
\textit{(Источник: Хливненко, Пятакович, 2023, с. 16--28; Хайкин, 2020, гл. 4.)}

\paragraph{Место нейросетей в классификации параллельных структур.}
С точки зрения классификации Флинна нейросети можно рассматривать как SIMD- или MIMD-структуры в зависимости от способа реализации и режима работы. Аппаратные нейропроцессоры, выполняющие одинаковые операции над множеством нейронов, тяготеют к SIMD; программные реализации на многоядерных процессорах или кластерах используют MIMD-параллелизм. Нейросети представляют собой проблемно-ориентированную параллельную структуру, эффективность которой определяется степенью параллелизма вычислений и возможностями аппаратной поддержки. В учебнике Орлова и Цилькера нейроинформатика отнесена к перспективным направлениям, которые будут определять развитие вычислительной техники в XXI веке.
\textit{(Источник: Хливненко, Пятакович, 2023, с. 3--5; Хайкин, 2020, гл. 1.)}

\paragraph{Применение нейросетей.}
Нейросети находят широкое применение в различных областях: распознавание образов и речи, обработка естественного языка (машинный перевод, анализ тональности), компьютерное зрение, прогнозирование временных рядов, управление роботами, медицинская диагностика, финансовый анализ и многие другие. Развитие аппаратных ускорителей (GPU, TPU, нейроморфных чипов) делает нейросетевые вычисления всё более эффективными. В пособии Хливненко и Пятакович особое внимание уделено практическим аспектам нейросетевого моделирования: визуализации внутреннего состояния обученной сети, решению задач классификации, категоризации, прогнозирования, восстановления зашумлённой информации, а также методам комбинирования градиентных и стохастических алгоритмов обучения для повышения эффективности решения практических задач.
\textit{(Источник: Хливненко, Пятакович, 2023, аннотация; Хайкин, 2020, аннотация.)}

\subsubsection{Краткое резюме подраздела}
Проблемно-ориентированные параллельные структуры включают матричные ВС, систолические структуры и нейросети. Матричные ВС реализуют SIMD-параллелизм на уровне массива процессорных элементов и эффективны для задач, хорошо распараллеливаемых по данным (обработка изображений, сеточные задачи, матричные вычисления). Систолические структуры обеспечивают конвейерно-параллельную обработку потоков данных с многократным использованием каждого элемента; они эффективны при матричных вычислениях, обработке сигналов и сортировке. Нейросети представляют собой параллельные структуры из простых нейронов с взвешенными связями, способные к обучению и обобщению; они относятся к перспективным направлениям развития вычислительной техники. Каждый класс структур имеет свои достоинства и недостатки, определяющие области его применения.
\textit{(Обобщение по: Орлов, Цилькер, гл. 13, с. 526--556; заключение, с. 598.)}

\end{document}